En el laboratori es reprodueix l'experiència d'Oersted. Es fa circular un corrent continu en el sentit que s'indica a la figura (eix $z$ positiu).

\figura[0.3]{fig1.png}

\begin{apartat}{1,25}
Representa esquemàticament l'orientació d'una brúixola situada als punts A i B que estan situats damunt l'eix y. Indica clarament la posició dels pols de la brúixola. Justifica la direcció i sentit de la brúixola.
\end{apartat}

\begin{apartat}{1,25}
Si s'inverteix la polaritat de la font de corrent continu que alimenta el fil, com variarà l'orientació de la brúixola als punts $A$ i $B$? Justifica la resposta. Quina informació ens aporta aquesta experiència sobre les fonts del camp magnètic?
\end{apartat}
